\section{AI for Math：形式化推理与 LLM Agent}

\subsection{纯神经方法的局限}

近年来 LLM 在数学竞赛中表现出色（如 O1 在 AMC 2024 中的显著提升、AlphaProof 在 IMO 中获得银牌级表现），但纯神经方法存在根本性弱点：

\begin{warningbox}{纯语言模型做数学的两个根本问题}
\begin{enumerate}
    \item \textbf{数据稀缺}：严格数学推理需要大量高质量推理轨迹或高质量奖励函数，超出竞赛/高中水平后难以仅靠人类数据获取
    \item \textbf{自然语言推理难以验证}：语言推理的正确性无法被形式化检验，在边界情况关键的实际系统中尤为危险
\end{enumerate}
\end{warningbox}

\subsection{形式化表示：替代方案}

\begin{importantbox}{形式化方法的核心思路}
\begin{enumerate}
    \item \textbf{自动形式化}（Autoformalization）：将非形式化的数学陈述转换为 Lean、Coq、Isabelle 等形式化语言
    \item \textbf{神经证明器}（Neural Prover）：在形式化环境中搜索证明策略（tactics）序列
    \item \textbf{验证}：形式化证明可被证明助手自动检验——这解决了验证难题
    \item \textbf{合成数据}：可大量生成候选证明并立即获得正确性反馈，将证明问题转化为类游戏任务
\end{enumerate}
\end{importantbox}

\subsubsection{Lean 证明助手工作原理}

Lean 是一个状态机：
\begin{itemize}
    \item \textbf{状态}（State）：当前需要证明的目标（goals）和上下文假设
    \item \textbf{动作}（Tactic）：简化规则，应用后改变证明状态
    \item \textbf{目标}：找到一系列 tactics，使所有 goals 被消解，到达 QED 状态
\end{itemize}

例如证明"若 $x$ 为偶数，则 $x^2$ 为偶数"：初始状态包含目标 $x^2 \bmod 2 = 0$，通过 \texttt{intro}、\texttt{rw} 等 tactics 逐步简化，最终达到 $0 = 0$ 即 QED。

\subsection{Copra：基于 LLM Agent 的定理证明}

\begin{importantbox}{Copra 系统架构}
Copra 是一种基于上下文学习（In-Context Learning）的定理证明 Agent：
\begin{enumerate}
    \item 使用前沿 LLM（如 GPT-4）预测下一步 tactic
    \item 将证明状态和搜索历史作为 Prompt 提供给 LLM
    \item 执行预测的 tactic，获取新状态或错误反馈
    \item 将错误信息反馈到 Prompt 中，让 LLM 修正策略
    \item 外层使用深度优先搜索（DFS）组织整个证明过程
\end{enumerate}
\end{importantbox}

Copra 的核心优势：
\begin{itemize}
    \item \textbf{无需额外训练}：直接利用预训练 LLM 的能力
    \item \textbf{即时受益于 LLM 进步}：从 GPT-4 到 O1/O3 的升级直接提升性能
    \item \textbf{融合自然语言与形式语言推理}
    \item \textbf{无需训练语料}：适用于全新的研究问题
    \item \textbf{证明助手反馈提供接地}（grounding）：防止幻觉
\end{itemize}

\subsection{层次化证明：抽象的力量}

\begin{importantbox}{LLM 驱动的层次化证明分解}
对于复杂定理（如 IMO 问题），Copra 可扩展为层次化求解：
\begin{enumerate}
    \item 让 LLM 生成非形式化的证明计划（informal proof plan）
    \item 基于计划将定理分解为多个子目标（sub-goals）的形式化表述
    \item 用 Copra 逐个证明子目标
    \item 将已证明的子目标放入上下文，证明完整定理
\end{enumerate}
关键洞察：LLM 不仅可以预测低级证明步骤，还可以进行\textbf{高级抽象推理}——这一切通过 Prompt 工程即可实现，无需额外训练。
\end{importantbox}

以 IMO 题目"证明 $\frac{21n+4}{14n+3}$ 对所有自然数 $n$ 不可约"为例：O3 将问题分解为关于 GCD 的三个子引理，Copra 分别证明后，再组合证明完整定理。

\subsection{应用：编译器形式化验证}

\begin{knowledgebox}{形式化验证的实际价值}
形式化验证可以数学地证明系统的正确性、安全性和性能。美国国防部曾使用形式化验证的 Linux 内核（seL4）部署在无人机上，红队6周内无法攻破。

历史上形式化验证极其昂贵，但 AI for Math 的兴起正在\textbf{根本性改变成本-收益方程}。
\end{knowledgebox}

讲座演示了一个简单的算术表达式编译器验证案例（Coq语言）：
\begin{itemize}
    \item 源语言：递归定义的算术表达式（加法、乘法）
    \item 目标语言：基于栈机的指令列表
    \item 编译器正确性定理：编译后执行结果 = 源语言求值结果
    \item Copra 无法一次性证明，但通过 LLM 发明辅助引理（单条指令编译正确性），再组合完成完整证明
\end{itemize}

\subsection{本章小结}

\begin{itemize}
    \item LLM Agent 用于数学发现极具前景，上下文学习方法已展现强大能力
    \item 证明助手的反馈对于接地（grounding）至关重要
    \item 高级 Agent 设计可组合形式化、层次化规划和低级证明生成
    \item 目前仍需 Agent 架构，纯训练方法尚未达到足够水平
\end{itemize}

